\uuid{cxS2}
\exo7id{515}
\titre{exo7 515}
\auteur{cousquer}
\organisation{exo7}
\datecreate{2003-10-01}
\isIndication{false}
\isCorrection{false}
\chapitre{Suite}
\sousChapitre{Convergence}
\cpge{mpsi-mp2i.nombres-reels-suites.d, pcsi.nombres-reels-suites.c, ptsi.nombres-reels-suites.c, mpsi-mp2i.nombres-reels-suites.i, pcsi.nombres-reels-suites.g, ptsi.nombres-reels-suites.g, bcpst1.suites-reelles, bcpst1.suites-reelles-usuelles}
\module{Analyse}
\niveau{L1}
\difficulte{}

\contenu{
\question{
Étant donné $k\in\mathbb{R}_+$, que peut-on dire d'une suite 
$(u_n)$ qui vérifie
$\lim_{n\to\infty} \frac{u_{n+1}}{u_n} = k$ ?
\textit{Application} : Étudier $u_n=\frac{1\cdot 2 \cdots n}{1\cdot 4\cdots(3n-2)}$.
}
}
